\section{Introduction}

Modern language-model systems have multiple ways to answer knowledge-intensive questions. They can rely on parametric memory, retrieve external evidence, cache useful facts in external memory, apply temporary task-time updates, or consolidate stable updates more durably. Prior work usually optimizes only one of these levers at a time: retrieval selection and self-reflective retrieval policies \citep{asai2024selfrag,yan2024crag}, or editing systems that route knowledge into side memories or detachable patches \citep{wang2024wise,yu2023melo,zhang2024keStudy}. The central question in this project is different: when should one controller choose among all of these operations under one objective?

That question matters because different regimes reward different choices. On recurring factual queries, repeated retrieval may be wasteful once a reusable answer has been observed and stored, a setting directly motivated by PopQA-style factual access patterns \citep{mallen2022trust}. On updated knowledge, cached memory may become stale, retrieval may lag visible evidence, and temporary or durable adaptation may become preferable; this is the regime that freshness-oriented benchmarks such as FreshQA were designed to stress \citep{vu2024freshllms}. A controller that sees all of these options under one reward is therefore more expressive than a retrieval gate, a memory writer, or a standalone editor.

This paper makes two narrow, artifact-backed claims. First, on PopQA, the controller matches retrieval-level quality overall after a conservative non-recurring guardrail fix while reducing repeated retrieval on recurring queries. Second, on a deterministic bundled freshness/update benchmark, the controller improves updated-knowledge behavior over retrieval-only and memory-only baselines while exercising temporary adaptation, rare guarded consolidation, and rollback.

The contribution is a unified multi-timescale routing perspective backed by two complementary experiments:
\begin{itemize}[leftmargin=1.5em]
\item \textbf{PopQA}: reusable external memory and delayed-utility routing reduce repeated retrieval without sacrificing quality.
\item \textbf{Freshness}: temporary adaptation and guarded durable consolidation improve updated-knowledge behavior under causal evidence visibility.
\end{itemize}

